\documentclass[11pt]{article}

% ---------- Packages ----------
\usepackage[margin=1in]{geometry}
\usepackage{amsmath,amssymb,amsthm}
\usepackage{enumitem}
\usepackage{booktabs}
\usepackage{longtable}
\usepackage{array}
\usepackage{xcolor}
\usepackage{titlesec}
\usepackage{fancyhdr}
\usepackage{listings}
\usepackage[colorlinks=true,linkcolor=blue!50!black,citecolor=blue!50!black,urlcolor=blue!50!black]{hyperref}
\usepackage{parskip}
\setlength{\emergencystretch}{3em}

% ---------- Colors ----------
\definecolor{hubblue}{RGB}{31,73,125}
\definecolor{extgray}{RGB}{110,110,110}
\definecolor{codebg}{RGB}{245,245,245}

% ---------- Section styling ----------
\titleformat{\section}{\normalfont\Large\bfseries\color{hubblue}}{\thesection}{1em}{}
\titleformat{\subsection}{\normalfont\large\bfseries\color{hubblue}}{\thesubsection}{1em}{}
\titleformat{\subsubsection}{\normalfont\bfseries}{\thesubsubsection}{1em}{}

\pagestyle{fancy}
\fancyhf{}
\fancyhead[L]{\small\itshape Chapter-1 Computational Companion}
\fancyhead[R]{\small\itshape\rightmark}
\fancyfoot[C]{\thepage}
\renewcommand{\headrulewidth}{0.3pt}

% ---------- Code listing style ----------
\lstset{
  basicstyle=\ttfamily\small,
  backgroundcolor=\color{codebg},
  frame=single,
  rulecolor=\color{gray!40},
  breaklines=true,
  columns=fullflexible,
  showstringspaces=false,
  xleftmargin=8pt,
  xrightmargin=8pt,
  aboveskip=8pt,
  belowskip=8pt
}

% ---------- Custom commands ----------
\newcommand{\extnote}[1]{{\small\color{extgray}\emph{[External / not in Pavliotis: #1]}}}
\newcommand{\pav}[1]{{\small\color{hubblue}(Pavliotis, #1)}}
\newcommand{\Dstar}{D_{\star}}

% Unit metadata block
\newenvironment{unitmeta}{%
  \begin{description}[style=sameline,leftmargin=1.4cm,font=\bfseries\color{hubblue}]
}{%
  \end{description}
}

% Per-unit implementation-contract signature table
\newenvironment{contracttable}{%
  \par\smallskip\noindent{\bfseries\color{hubblue}Implementation contract.}\par\nopagebreak\smallskip
  \small
  \setlength{\extrarowheight}{2pt}%
  \noindent\begin{tabular}{@{}p{3.5cm}p{6.8cm}p{4.2cm}@{}}
  \toprule
  \textbf{Role} & \textbf{Call signature} & \textbf{Produces / returns} \\
  \midrule
}{%
  \bottomrule
  \end{tabular}\par\smallskip
}

\title{\textbf{A Computational Companion to Pavliotis, Chapter 1}\\[4pt]
\large Project Plan: Covariance Operators, Spectral Theory, and the Karhunen--Lo\`eve Expansion}
\author{Stefan}
\date{\today}

\begin{document}

\maketitle

\begin{abstract}
\noindent This document specifies a six-to-eight-week computational project in Julia that deepens understanding of Chapter~1 of Pavliotis, \emph{Stochastic Processes and Applications}, by implementing its core objects and verifying each implementation against an analytic prediction. The project is organized around a single hub --- the covariance operator $\mathcal{C}$ --- with two diagonalizing bases (Fourier/Bochner and Mercer/Karhunen--Lo\`eve), \emph{four} factorizations used for sampling (Cholesky, KL truncation, random Fourier series, and circulant embedding), and ergodicity as the mechanism that closes the loop between one simulated sample path and the analytic objects it is checked against. The plan deliberately excludes machine-learning and SciML content. At the working budget of \textbf{15\,h/week} ($6\times 2.5$\,h) the committed core is \textbf{Units 0--7}: the covariance/spectral spine (Units 0--4), Brownian motion as a scaling limit (Unit~5), fractional Brownian motion (Unit~6), and a forward-pointing SDE bridge to Chapter~3 (Unit~7). Each unit now carries an explicit \emph{implementation contract} (the \texttt{src/} functions it exercises, with signatures, and the data its \texttt{run.jl} must produce), an \emph{exact analytic target}, a \emph{pass/fail predicate}, a \emph{negative control} that shows the wrong thing failing, and a one-line \emph{README argument} --- since the READMEs, not the code, are the durable evidence of understanding. Statements drawn directly from the text are cited by section/theorem number; statements that are standard but external to Pavliotis are flagged explicitly. All citations are to the project's copy of the text, which is the \emph{November 11, 2015 draft} (numbering differs from the published edition), and were checked directly against it.
\end{abstract}

\tableofcontents

% =====================================================================
\section{The Through-Line}
% =====================================================================

\subsection{The hub: the covariance operator}

Every object studied in Chapter~1 is a fact about a single operator. For a mean-zero, second-order process $X_t$, $t\in[0,T]$, with continuous covariance $R(t,s) = \mathbb{E}(X_tX_s)$, define the integral operator
\begin{equation}
(\mathcal{C}f)(t) \;:=\; \int_0^T R(t,s)\,f(s)\,ds, \qquad f \in L^2[0,T].
\label{eq:covop}
\end{equation}
\pav{eqn.~1.33} This is self-adjoint, nonnegative, and (for continuous $R$) compact and Hilbert--Schmidt \pav{Exercises 21--22}. Everything downstream is one of two operations performed on $\mathcal{C}$.

\subsection{Two diagonalizations}

\paragraph{Fourier / Bochner (gated by stationarity).}
If $X_t$ is weakly stationary, $R(t,s) = R(t-s)$ is a function of a single lag, and Bochner's theorem applies:
\begin{equation}
R(\tau) \;=\; \int_{\mathbb{R}} e^{i\omega\tau}\,\rho(d\omega),
\label{eq:bochner}
\end{equation}
for a unique nonnegative measure $\rho$ on $\mathbb{R}$ with $\rho(\mathbb{R}) = R(0)$ \pav{Thm.~1.14}. When $\rho$ is absolutely continuous, $\rho(d\omega) = S(\omega)\,d\omega$ with
\begin{equation}
S(\omega) = \frac{1}{2\pi}\int_{-\infty}^{\infty} e^{-i\omega\tau}R(\tau)\,d\tau,
\qquad
R(\tau) = \int_{-\infty}^{\infty} e^{i\omega\tau}S(\omega)\,d\omega.
\label{eq:spectraldensity}
\end{equation}
\pav{eqns.~1.7--1.8} $S(\omega)$ is the \emph{spectral density}; it lives on the frequency axis, not on the sample space --- a distinction worth stating explicitly since both objects are conventionally denoted $\omega$. Equations~\eqref{eq:spectraldensity} form a Fourier transform pair and invert by the Fourier inversion theorem, independently of any symmetry of $R$. Realness and evenness of $R$ add only that $S$ is itself real and even --- so the sign of $i$ in the exponent is immaterial (one may replace $e^{\pm i\omega\tau}$ by $\cos\omega\tau$); the $1/2\pi$ is placed on the $S$-side purely by convention. With this convention $\int_{\mathbb{R}} S(\omega)\,d\omega = R(0)$ (total variance), which fixes the normalization every estimator in Unit~1 must reproduce.

\paragraph{Mercer / Karhunen--Lo\`eve (general, needs a compact domain).}
Independently of stationarity, the spectral theorem for compact self-adjoint operators gives a countable orthonormal eigenbasis $\{e_k\}$ of $\mathcal{C}$ with eigenvalues $\lambda_k \geq 0$,
\begin{equation}
\mathcal{C}e_k = \lambda_k e_k, \qquad R(t,s) = \sum_{k=1}^{\infty}\lambda_k\,e_k(t)\,e_k(s),
\label{eq:kl-eigen}
\end{equation}
\pav{eqn.~1.32; Mercer's theorem, stated inline (unnumbered) in \S1.5} and the Karhunen--Lo\`eve expansion
\begin{equation}
X_t = \sum_{k=1}^{\infty} \xi_k\, e_k(t), \qquad \xi_k = \int_0^T X_t\, e_k(t)\,dt,\quad \mathbb{E}(\xi_k\xi_\ell) = \lambda_k\delta_{k\ell},
\label{eq:kl}
\end{equation}
\pav{Thm.~1.28 (Karhunen--Lo\`eve)} converging in $L^2$, uniformly in $t$. If $X_t$ is Gaussian, the $\xi_k$ are independent (uncorrelated \emph{and} jointly Gaussian), and
\begin{equation}
X_t = \sum_{k=1}^{\infty} \sqrt{\lambda_k}\,\xi_k\, e_k(t), \qquad \xi_k \overset{\text{iid}}{\sim} N(0,1).
\label{eq:kl-gaussian}
\end{equation}
\pav{eqn.~1.36} The eigenvalues are simultaneously (i) eigenvalues of the covariance operator and (ii) variances of the independent stochastic coefficients in the expansion --- one fact, two readings.

\paragraph{When the two diagonalizations coincide.}
The Fourier basis diagonalizes $\mathcal{C}$ exactly when $\mathcal{C}$ is a convolution operator, which requires the domain to be a group under addition and the kernel to be translation-invariant. The interval $[0,T]$ is not a group; the circle (torus) is. On the torus, $\mathcal{C}$ is a genuine convolution, its eigenfunctions are forced to be the characters $e^{ikt}$, and
\begin{equation}
\lambda_k = \hat R(k) \qquad \text{(exactly)},
\label{eq:torus-coincidence}
\end{equation}
the $k$-th eigenvalue equals the $k$-th Fourier coefficient of $R$ --- the KL sequence and the (discrete, atomic) spectral measure are the same object. On $[0,T]$ the boundary breaks translation invariance: the KL eigenfunctions of Brownian motion, for instance, are Sturm--Liouville sine modes with eigenvalues $\lambda_n \propto (2n-1)^{-2}$ \pav{Example 1.29}, not Fourier modes, and \eqref{eq:torus-coincidence} fails. \extnote{the general asymptotic recovery $\lambda_k \to S(\omega_k)$ as $T\to\infty$ is a Toeplitz-operator fact (Grenander--Szeg\H{o}), not covered in Pavliotis. This asymptotic --- not an exact $[0,T]$ identity --- is precisely the object of the Unit-3 cross-check.}

\subsection{One factorization, four routes}

Sampling a Gaussian process is applying a square root of $\mathcal{C}$ to white noise. Four square roots give four sampling algorithms, all reproducing the same \emph{law}:
\begin{itemize}[leftmargin=1.4cm]
  \item \textbf{Cholesky (general):} $\Sigma = LL^\top$ on a grid, $X = Lz$, $z\sim N(0,I)$.
  \item \textbf{KL truncation (general):} \eqref{eq:kl-gaussian} truncated at $K$ terms.
  \item \textbf{Random Fourier series, Gaussian amplitudes (stationary only):} on a \emph{one-sided} frequency grid $\omega_k>0$,
  $X(t) = \sum_k \big[a_k\cos(\omega_k t) + b_k\sin(\omega_k t)\big]$ with $a_k,b_k \overset{\text{iid}}{\sim} N\!\big(0,\,2\,S(\omega_k)\,\Delta\omega\big)$. The factor $2$ folds the two-sided spectrum onto $\omega>0$: $\mathrm{Cov}(\tau)=\sum_k 2S(\omega_k)\Delta\omega\cos(\omega_k\tau)\to R(\tau)$. \emph{Gaussian} amplitudes make this a genuine square root of $\mathcal{C}$ applied to white noise, so the finite-$K$ sample is exactly Gaussian; a random-\emph{phase} variant (fixed amplitude $\sqrt{4S(\omega_k)\Delta\omega}$, uniform phase) reproduces the covariance but \emph{not} the Gaussian law at finite $K$ (Gaussian only in the limit, by the CLT), and is therefore not a square root.
  \item \textbf{Circulant embedding (stationary only):} FFT-diagonalization of the circularly-extended Toeplitz covariance (Wood--Chan / Dietrich--Newsam); the circulant eigenvalues are the DFT of the covariance sequence --- i.e.\ the \emph{discrete} spectral density --- so this is precisely the Bochner square root made computational, on the grid rather than on the line. Exact and $O(m\log m)$ when the embedding is nonnegative definite.
\end{itemize}
Cholesky and KL are general; the random-Fourier and circulant routes require stationarity. Each is a different basis in which to take the square root of the same operator: Cholesky is the triangular factor, KL is the eigenbasis, and random-Fourier / circulant embedding are the Bochner square roots on the continuous frequency axis and the discrete FFT grid respectively.

\subsection{Stationarity as gate, ergodicity as loop-closer}

Stationarity is a \emph{precondition} for the Fourier vertex: Brownian motion is not stationary (its covariance $\min(t,s)$ is not a function of $t-s$), so it has no spectral density; it possesses only the Mercer/KL vertex. Ergodicity is the mechanism that licenses every theoretical check in this project: for a stationary process with $C\in L^1(0,\infty)$,
\begin{equation}
\lim_{T\to\infty}\mathbb{E}\left|\frac{1}{T}\int_0^T X_s\,ds - \mu\right|^2 = 0,
\label{eq:ergodic}
\end{equation}
\pav{Prop.~1.16, eqn.~1.12} with finite-$T$ rate (for $C\geq 0$)
\begin{equation}
\mathbb{E}\left|\frac{1}{T}\int_0^T (X_s-\mu)\,ds\right|^2 = \frac{2}{T^2}\int_0^T (T-u)\,C(u)\,du \;\leq\; \frac{2}{T}\int_0^\infty C(u)\,du \;=\; \frac{2\Dstar}{T},
\label{eq:ergodic-rate}
\end{equation}
\pav{eqn.~1.13, Lemma 1.17} where the first equality is general (Lemma~1.17) but the inequality uses $C\geq0$ (for sign-changing $C$ the honest bound is $\tfrac{2}{T}\int_0^\infty|C|$, and $\Dstar=\int_0^\infty C$ coincides with it only in the sign-definite case; the asymptotic form $\mathrm{Var}\sim 2\Dstar/T$ holds regardless of sign given integrability). Here we introduce the \textbf{Green--Kubo (transport) coefficient}
\begin{equation}
\Dstar \;:=\; \int_0^\infty C(t)\,dt,
\qquad
\mathbb{E}\Big(\int_0^t X_s\,ds\Big)^2 \approx 2\Dstar t \ \ (t \gg 1).
\label{eq:greenkubo}
\end{equation}
\pav{eqns.~1.14--1.15} \textbf{Notation caution (carry $\alpha$).} Pavliotis denotes both the transport coefficient $\int_0^\infty C$ (eqn.~1.14) \emph{and} the noise strength in the OU kernel via $C(0)=D/\alpha$ (Example~1.15) by the single symbol $D$. Example~1.18 makes the collision explicit: there Pavliotis computes the diffusion coefficient of the exponential-correlation process as $\int_0^\infty C = C(0)\,\tau_{\mathrm{cor}} = D/\alpha^2$, so for the OU kernel the two quantities differ, $\Dstar = \int_0^\infty C = D/\alpha^2$, coinciding only at $\alpha=1$. (Eqn.~1.11 merely sets $\alpha=D=1$ for a specific covariance computation, incidentally collapsing the distinction --- it is not the locus of the conflation.) We keep $D$ for the noise strength and $\Dstar$ for the transport coefficient throughout, so every rate constant carries its explicit $\alpha$-dependence. This is what justifies estimating a covariance or a spectral density from a \emph{single} simulated path and comparing it to a closed-form target --- the core verification methodology used throughout Section~2.

\subsection{Summary}

\begin{center}
\begin{tabular}{@{}p{3.1cm}p{4.3cm}p{4.3cm}@{}}
\toprule
& \textbf{Fourier / Bochner} & \textbf{Mercer / KL} \\
\midrule
Requires & weak stationarity & compact domain (continuous $\mathcal{C}$) \\
Object & spectral density $S(\omega)$ & eigenpairs $(\lambda_k,e_k)$ \\
Lives on & frequency axis $\mathbb{R}$ & discrete index $k\in\mathbb{N}$ \\
Exact coincidence & \multicolumn{2}{c}{on the torus only: $\lambda_k=\hat R(k)$} \\
\bottomrule
\end{tabular}
\end{center}

% =====================================================================
\section{Units}
% =====================================================================

\textbf{Committed core: Units 0--7} at the working budget of $15$\,h/week ($6\times 2.5$\,h) against a $6$--$8$ week horizon, i.e.\ $90$--$120$\,h. Summed per-unit estimates below total $\approx 52$--$64$ core hours; the remaining budget is deliberate headroom for per-unit spec-writing and review of all generated code, plus a debugging buffer concentrated on the units most likely to overrun for someone newer to stochastic numerics: \textbf{Units~1, 2, 6, and 7}. The risk in Units~1, 6, and 7 is spectral/FFT normalization (one- vs two-sided, angular vs ordinary frequency, unnormalized FFTW, window power); the risk in Unit~2 is the Nystr\"om symmetrization and the eigenvalue noise floor. One further hidden time-sink to anticipate: the Unit-5 marginal-convergence rate is \emph{law-dependent} (see Unit~5), so a fitted slope that is not $n^{-1/2}$ can be correct rather than a bug. Everything through Unit~6 is firm; the SDE bridge (Unit~7) is a labelled forward pointer to Chapter~3 and is the single cut line if the buffer is ever consumed --- nothing in Units 0--6 depends on it. This protects depth over breadth per the stated priority: the extra budget over the original $7.5$\,h/week plan is spent on \emph{more of Chapter~1 done thoroughly} (the scaling-limit construction, fBm, the truncation-cost study, circulant embedding), not on rushing ahead.

Each unit below adds five implementation-oriented fields to the narrative ones: \textbf{Analytic target} (the exact closed-form object the check compares against), \textbf{Pass/fail} (the numerical gate as a checkable predicate), \textbf{Negative control} (where the \emph{wrong} thing is shown to fail), \textbf{README} (the one-line argument its \texttt{experiments/NN\_*/README.md} must make), and \textbf{Missteps} (the most probable errors, each with a one-line guard). A signature table then names the \texttt{src/} surface the unit exercises and the data its \texttt{run.jl} must emit.

% ---------------------------------------------------------------
\subsection*{Unit 0 --- Covariance Core and Gaussian Sampling via Cholesky}
\addcontentsline{toc}{subsection}{Unit 0 --- Covariance Core and Gaussian Sampling via Cholesky}
\begin{unitmeta}
\item[Concept] Gaussian process determined by mean and covariance \pav{App.~B.5}; Cholesky as a square root, $\Sigma = LL^\top$.
\item[Build] A \texttt{GaussianProcess} type over a kernel $R(t,s)$; grid assembly of $\Sigma$; Cholesky factorization with a small jitter (nugget) $\varepsilon I$ for numerical positive-definiteness; path sampler; the empirical-covariance estimator and its convergence study.
\item[Check] Empirical covariance of $N$ paths against the analytic $R$, measured in Frobenius norm as a function of $N$.
\item[Analytic target] $\|\widehat\Sigma_N - \Sigma\|_F \propto N^{-1/2}$, where $\Sigma_{ij}=R(t_i,t_j)$; the intercept is irrelevant, the exponent is the object.
\item[Pass/fail] Fit $\log\|\widehat\Sigma_N-\Sigma\|_F$ against $\log N$ over $N\in[10^2,10^4]$; gate $|\widehat{\text{slope}}+\tfrac12| < 2.5\,\mathrm{SE}(\widehat{\text{slope}})$ (stochastic-slope gate).
\item[Negative control] Set $\texttt{jitter}=0$ and factor a fine-grid exponential-kernel $\Sigma$: \texttt{cholesky} throws \texttt{PosDefException} because $\Sigma$ is numerically indefinite. The throw disappears once $\varepsilon\gtrsim10^{-10}$ --- so the nugget is shown to be load-bearing, not cosmetic.
\item[Artifact] Sample paths; \emph{error-vs-$N$ log--log plot} with the fitted $-\tfrac12$ slope annotated against the reference line --- the collapse of the fit onto slope $-\tfrac12$ is the feature that demonstrates success.
\item[README] Concept: a Gaussian process is fixed by mean and covariance, and sampling is applying a square root of $\Sigma$. Why this check: the empirical covariance must converge to the analytic $R$ at the Monte-Carlo rate, and the \emph{slope} certifies both the sampler and the estimator at once. What it proves: the atomic draw underlying every later unit is correct.
\item[Missteps] (i) \emph{Jitter too small} $\Rightarrow$ indefinite throw; \emph{too large} $\Rightarrow$ biased covariance. Guard: report $\varepsilon$ and confirm $\|\widehat\Sigma_N-\Sigma\|$ is not floored by $\varepsilon$. (ii) \emph{Path-matrix orientation:} \texttt{empirical\_cov} expects paths as columns ($n_{\text{grid}}\times N$); a transpose silently estimates the wrong matrix. Guard: assert \texttt{size(paths,1)==length(t\_grid)}. (iii) \emph{Global RNG:} use \texttt{StableRNG}, not \texttt{randn()} on the default stream.
\item[Purpose] The atomic operation underlying every later unit: how a Gaussian random function is drawn at all.
\item[Effort] $\approx 6$--$8$ h (includes repository skeleton).
\item[Prerequisite] $X = Lz,\ z\sim N(0,I) \Rightarrow \operatorname{Cov}(X) = LL^\top$ --- standard linear algebra, already familiar; the only new content is that a Gaussian \emph{process} is a consistent family of such draws. \pav{App.~B.5}
\end{unitmeta}

\begin{contracttable}
\texttt{gaussianprocess.jl} & \texttt{GaussianProcess(kernel; meanfn=t->0.0)} & GP over $R(t,s)$ \\
 & \texttt{assemble\_cov(gp, t\_grid)} & \texttt{Symmetric} $\Sigma_{ij}=R(t_i,t_j)$ \\
 & \texttt{assemble\_mean(gp, t\_grid)} & mean vector $m(t_i)$ \\
 & \texttt{empirical\_cov(paths)} & grid$\times$grid sample covariance; \texttt{paths} is $n_{\text{grid}}\times N$ \\
\texttt{sampling.jl} & \texttt{sample\_cholesky(Sigma, rng; jitter=1e-10)} & one path $X=Lz$ \\
\end{contracttable}
\noindent\textbf{\texttt{experiments/00\_covariance\_core/run.jl}} produces: (i) an $n_{\text{grid}}\times N$ path matrix and a sample-paths figure; (ii) the table $(N,\ \|\widehat\Sigma_N-\Sigma\|_F)$ over the $N$-ladder and the log--log error plot; (iii) the recorded jitter and RNG seed.

% ---------------------------------------------------------------
\subsection*{Unit 1 --- Stationarity Gate: Spectral Density vs.\ Bochner (with Circulant Embedding)}
\addcontentsline{toc}{subsection}{Unit 1 --- Stationarity Gate: Spectral Density vs.\ Bochner}
\begin{unitmeta}
\item[Concept] Weak stationarity \pav{\S1.2}; Bochner's theorem \eqref{eq:bochner}; the Ornstein--Uhlenbeck\slash exponential kernel and its Lorentzian spectrum,
\begin{equation}
R(\tau) = \frac{D}{\alpha}e^{-\alpha|\tau|} \quad\Longrightarrow\quad S(\omega) = \frac{D}{\pi}\,\frac{1}{\omega^2+\alpha^2}.
\label{eq:lorentzian}
\end{equation}
\pav{Example 1.15}
\item[Build] Stationary GP with kernel \eqref{eq:lorentzian}; estimate $S$ via an averaged (Welch) periodogram and via FFT of the estimated autocorrelation; retain the \emph{raw} periodogram as a foil; \textbf{implement circulant embedding} (Route~4 in \S1.3) as the exact FFT-based stationary sampler --- the same Bochner/FFT machinery running in the sampling direction.
\item[Check] Estimated $S(\omega)$ against the Lorentzian; the total-variance normalization; and circulant samples against the analytic covariance.
\item[Analytic target] $S(\omega)=\dfrac{D}{\pi(\omega^2+\alpha^2)}$ with $\displaystyle\int_{\mathbb{R}}S\,d\omega = D/\alpha = R(0)$ under the $1/2\pi$-on-$S$ convention of \eqref{eq:spectraldensity}.
\item[Pass/fail] (a) $\big|\int\widehat S\,d\omega\,/\,R(0) - 1\big| < 0.05$; (b) relative $L^2$ error of $\widehat S$ vs.\ the Lorentzian over the resolved band $|\omega|\le\omega_{\max}$ below a few times the Welch estimator's own SE (documented in the README); (c) circulant covariance error $\|\widehat\Sigma-\Sigma\|/\|\Sigma\| < 10^{-10}$ (floating-point scale).
\item[Negative control] Two, both plotted. (i) Overlay \texttt{raw\_periodogram} on \texttt{welch\_psd}: the raw estimate stays jagged and its variance does \emph{not} shrink with record length at fixed $\omega$, while Welch converges --- inconsistency made visible. (ii) Drop the $1/2\pi$ (engineering convention): $\int\widehat S\,d\omega\,/\,R(0)$ lands on exactly $2\pi$ instead of $1$ --- the spurious factor made to appear on purpose.
\item[Artifact] Estimated vs.\ analytic $S(\omega)$, log--log, showing the Lorentzian tail with the resolved band shaded; a raw-vs-Welch overlay; a circulant-embedding covariance-error panel.
\item[README] Concept: stationarity moves all second-order structure onto the frequency axis (Bochner). Why this check: the estimate must both match the Lorentzian \emph{and} integrate to the variance under a pinned convention, and circulant embedding is the same machinery run as a sampler. What it proves: the Fourier vertex is correctly implemented in both directions, and the two controls exhibit exactly what a wrong estimator and a wrong normalization look like. (This README carries the $2\pi$ lesson for the whole repository.)
\item[Missteps] (i) \emph{One- vs two-sided spectrum} --- forgetting the interior-bin factor $2$ halves $\int\widehat S$. (ii) \emph{Angular vs ordinary frequency} --- mixing $\omega$ and $f=\omega/2\pi$ injects a $2\pi$. (iii) \emph{FFTW is unnormalized} --- restore $\Delta t$ and $1/N$ explicitly. (iv) \emph{Welch window power} --- divide by $\sum_n w_n^2$, not $N$. (v) \emph{Circulant scaling} --- validate the exact FFT normalization against the analytic covariance here, once. Guard for all: pin the normalization so the discrete $\int\widehat S\,d\omega\to R(0)$ and unit-test it.
\item[Purpose] The Fourier vertex: stationarity is what allows second-order structure to be described entirely on the frequency axis --- both as an object to estimate (periodogram) and as an exact synthesis route (circulant embedding).
\item[Effort] $\approx 8$--$9$ h (includes circulant embedding). \emph{Highest overrun risk in the project} --- the circulant scaling validation and the two controls are where the hours go.
\item[Prerequisite] \extnote{the raw periodogram is a statistically \emph{inconsistent} estimator of $S$ and must be smoothed or averaged; this is a standard time-series fact (e.g.\ Percival \& Walden, or Brockwell \& Davis), not covered in Pavliotis.} The point genuinely new relative to deterministic FFT experience: the estimator itself is random and biased, so familiarity with the FFT does not by itself deliver a consistent density estimate.
\end{unitmeta}

\begin{contracttable}
\texttt{kernels.jl} & \texttt{exponential\_kernel(t, s; D=1.0, alpha=1.0)} & OU/exponential covariance \\
\texttt{spectral.jl} & \texttt{welch\_psd(x, dt; nseg, noverlap, window=:hann, onesided=true)} & consistent $(\omega,\widehat S)$ \\
 & \texttt{raw\_periodogram(x, dt; onesided=true)} & inconsistent $(\omega,I)$ (control) \\
 & \texttt{bochner\_forward(r, dt; onesided=true)} & $(\omega,S)$ from a covariance sequence \\
 & \texttt{spectral\_variance(omega, Shat)} & trapezoid $\int\widehat S\,d\omega$ \\
\texttt{sampling.jl} & \texttt{sample\_circulant\_embedding(r, rng)} & exact stationary path \\
\end{contracttable}
\noindent\textbf{\texttt{experiments/01\_spectral\_bochner/run.jl}} produces: (i) $(\omega,\widehat S,S_{\text{analytic}})$ and the log--log density plot; (ii) the scalar \texttt{spectral\_variance/R(0)} (target $\approx1$); (iii) $\|\widehat\Sigma-\Sigma\|/\|\Sigma\|$ for circulant samples; (iv) the raw-vs-Welch and dropped-$1/2\pi$ control panels.

% ---------------------------------------------------------------
\subsection*{Unit 2 --- Karhunen--Lo\`eve via Quadrature Eigenproblem (Torus Contrast + Truncation Cost)}
\addcontentsline{toc}{subsection}{Unit 2 --- Karhunen--Lo\`eve via Quadrature Eigenproblem}
\begin{unitmeta}
\item[Concept] KL expansion \pav{\S1.5, eqns.~1.32--1.36, Thm.~1.28}; Mercer's theorem \pav{inline, \S1.5}; closed-form Brownian-motion KL basis on $[0,1]$,
\begin{equation}
\lambda_n = \left(\frac{2}{(2n-1)\pi}\right)^{2}, \qquad e_n(t) = \sqrt{2}\,\sin\!\left(\tfrac{1}{2}(2n-1)\pi t\right).
\label{eq:bm-kl}
\end{equation}
\pav{Example 1.29, eqn.~1.37} (Pavliotis fixes $T=1$ in \S1.5; on a general $[0,T]$ these scale as $\lambda_n \propto T^2$, and the trace target below is $T$-dependent.)
\item[Build] Nystr\"om discretization of the eigenproblem $\int_0^T R(t,s)e(s)\,ds = \lambda\, e(t)$ by quadrature, \textbf{symmetrized} before the eigensolve (see Missteps), reducing it to a symmetric matrix eigenproblem --- directly analogous to boundary-integral-equation discretization on a positive semi-definite kernel; sample by KL truncation.
\item[Check] Numerical $(\lambda_k,e_k)$ against the closed form; the trace identity as an assembly check; the torus coincidence and its failure on $[0,T]$; the truncation cost/accuracy trade-off.
\item[Analytic target] Eigenpairs \eqref{eq:bm-kl} ($\lambda_n\propto T^2$ on $[0,T]$); trace $\sum_n\lambda_n = \int_0^T R(t,t)\,dt$ (general form; $=T^2/2$ for Brownian motion, \emph{not} $T\,R(0)=0$), with the stationary specialization $\sum_n\lambda_n=T\,R(0)$ \pav{Exercise 23} checked on a stationary kernel; torus coincidence $\lambda_k=\hat R(k)$ exactly.
\item[Pass/fail] Eigenvalue error and eigenfunction $L^2$-error vs.\ $k$, gated \emph{only over the resolved range} above the discretization floor; the deterministic $\lambda_k\sim k^{-2}$ decay slope fitted within that window and gated against $-2$ with a fixed absolute margin ($\pm0.15$; \emph{not} an SE-based gate --- the residuals are misspecification, not sampling noise). Trace: $|\sum\lambda_{\text{num}}-\int_0^T R(t,t)\,dt|/|\!\int_0^T R(t,t)\,dt| < 10^{-3}$. Torus: $\max_k|\lambda_k-\hat R(k)| < 10^{-8}$ on the circle.
\item[Negative control] Two, both latent in the design and now labelled. (i) Run the torus coincidence $\lambda_k=\hat R(k)$ on $[0,T]$: it \emph{fails}, because the eigenfunctions are Sturm--Liouville sines, not characters. (ii) Extend the eigenvalue plot \emph{below} the discretization floor: the $k^{-2}$ slope flattens --- the resolution limit is made visible rather than hidden.
\item[Artifact] First eigenfunctions; eigenvalue decay $\lambda_k \sim k^{-2}$ with the resolved window shaded and the floor annotated; side-by-side torus-Fourier vs.\ interval-Sturm--Liouville mode shapes; truncation-error-vs-$K$ and KL-vs-Cholesky cost curves.
\item[README] Concept: KL is the adapted-basis diagonalization, and its eigenvalues double as the variances of the expansion coefficients (the classical antecedent of POD/PCA). Why this check: Brownian motion's closed-form basis certifies the solver, the trace pins assembly, and the torus contrast pins the Fourier$\leftrightarrow$KL relationship --- exact on the circle, broken on the interval. What it proves: the second diagonalization is correctly implemented and its exact relationship to the first is understood, with the resolved range stated so the slope claim is honest.
\item[Missteps] (i) \emph{Non-symmetric discretization (the key one):} $\int R(t,s)e(s)\,ds=\lambda e(t)$ discretizes to $(KW)e=\lambda e$ with $W=\operatorname{diag}(\text{weights})$, which is \emph{not} symmetric even though $R$ is --- its eigenvectors are not orthonormal and may come out complex. Guard: solve $W^{1/2}KW^{1/2}g=\lambda g$, then $e=W^{-1/2}g$. (ii) \emph{Eigenvector normalization:} normalize in the discrete $L^2$ inner product ($\sum_i w_i e_k(t_i)^2=1$), not the Euclidean one. (iii) \emph{Sign convention:} eigenvector signs are arbitrary; align sign before computing the eigenfunction $L^2$ error. (iv) \emph{Noise floor:} fit the decay slope only above it.
\item[Numerical caution] The small eigenvalues are computed to poor \emph{relative} accuracy: Nystr\"om resolves $\lambda_k$ only until it sinks into the discretization noise floor, after which the error-vs-$k$ curve plateaus and the $\lambda_k\sim k^{-2}$ decay flattens. This floor is a property of the discretization, not a modeling error; the slope checks are meaningful only over the $k$-range above it, and the experiment README states the resolved range explicitly.
\item[Purpose] The adapted-basis diagonalization: the optimal linear basis for representing a given random function (the classical antecedent of POD/PCA), together with the practical cost/accuracy trade-off of truncating it.
\item[Effort] $\approx 10$--$11$ h --- the deepest unit (includes the truncation-cost study); flagged for the debugging buffer (symmetrization + noise floor).
\item[Prerequisite] The spectral theorem for compact self-adjoint operators (already familiar from functional analysis coursework); the new content is that the kernel being diagonalized is specifically a \emph{covariance}, so its eigenvalues carry the further meaning of variances. \pav{\S1.5, Thm.~1.28 and the inline Mercer statement}
\end{unitmeta}

\noindent The trace identity used as the assembly/quadrature sanity check is, in its general form,
\begin{equation}
\sum_{n=1}^{\infty}\lambda_n \;=\; \operatorname{Tr}\mathcal{C} \;=\; \int_0^T R(t,t)\,dt
\label{eq:trace}
\end{equation}
which equals $T^2/2$ for Brownian motion (\emph{not} $T\,R(0)=0$); the stationary specialization $\sum_n\lambda_n=T\,R(0)$ \pav{Exercise 23} holds only when $R(t,t)\equiv R(0)$ and is checked separately on a stationary kernel.

\noindent\emph{Torus contrast, in full.} On the circle with a stationary kernel, verify $\lambda_k = \hat R(k)$ \emph{exactly} \eqref{eq:torus-coincidence} with Fourier eigenfunctions, and confirm this identity fails on $[0,T]$ where the eigenfunctions are the Sturm--Liouville sines of \eqref{eq:bm-kl}. (Pavliotis's Exercise~27 verifies $\lambda_k=\hat R(k)$ only for the \emph{degenerate}, rank-2 kernel $\cos 2\pi(t-s)$; the project's \texttt{periodic\_kernel} has a full positive spectrum $\hat R(k)=\tfrac{2\alpha\,(1-(-1)^k e^{-\alpha/2})}{\alpha^2+(2\pi k)^2}>0$, derived and verified in-house, so the torus coincidence here rests on that derivation rather than on Exercise~27.) \emph{Truncation cost/accuracy (Exercise~28):} quantify the KL truncation error as a function of the number of retained modes $K$ (how many terms are needed for accurate statistics of BM, the bridge, and OU) and compare the cost of KL sampling against the Cholesky route of Unit~0.

\begin{contracttable}
\texttt{kl.jl} & \texttt{quad\_nodes\_weights(T; n, rule=:trapezoid)} & \texttt{(nodes, weights)} \\
 & \texttt{nystrom\_eigen(R, nodes, weights; nev=n)} & \texttt{(lambdas, eigfuncs)}, symmetrized, sorted, $L^2$-normalized \\
 & \texttt{trace\_diag(R, nodes, weights)} & $\sum_i w_i R(t_i,t_i)\approx\int_0^T R(t,t)\,dt$ \\
 & \texttt{kl\_tail\_energy(lambdas, K)} & $\sum_{k>K}\lambda_k/\sum_k\lambda_k$ \\
\texttt{sampling.jl} & \texttt{sample\_kl(lambdas, eigenfuncs, rng)} & one KL-truncated path \\
\texttt{kernels.jl} & \texttt{periodic\_kernel(t, s; D=1.0, alpha=1.0)} & torus kernel (positive spectrum) \\
\end{contracttable}
{\sloppy
\noindent\textbf{\texttt{experiments/02\_kl\_quadrature/run.jl}} produces: (i) $(k,\lambda_{\text{num}},\lambda_{\text{analytic}})$ and $(k,\text{eigenfunction }L^2\text{-error})$ over the resolved range; (ii) $\sum\lambda-\texttt{trace\_diag}$ (general) and the stationary $T\,R(0)$ check; (iii) the torus $(k,\lambda_k,\hat R(k))$ exact-match table and the same identity failing on $[0,T]$; (iv) $(K,\ \texttt{kl\_tail\_energy})$ and the KL-vs-Cholesky cost table.\par}

% ---------------------------------------------------------------
\subsection*{Unit 3 --- Reconciliation: The Process Zoo Where the Machinery Agrees with Itself}
\addcontentsline{toc}{subsection}{Unit 3 --- Reconciliation (Process Zoo)}
\begin{unitmeta}
\item[Concept] Brownian motion, $R(t,s)=\min(t,s)$ \pav{eqn.~1.21}, and its invariances (rescaling, shifting, time reversal, inversion) \pav{Prop.~1.23}; Brownian bridge, $R(t,s) = \min(t,s) - ts$, with $B_t = (1-t)\,W\!\left(\frac{t}{1-t}\right)$ \pav{eqn.~1.26 / Exercise 6}; stationary Ornstein--Uhlenbeck, $R(\tau) = \frac{D}{\alpha}e^{-\alpha|\tau|}$ \pav{Example 1.15}. The unit's one distinct idea is \emph{reconciliation}: the sampling routes and the two diagonalizations are forced to agree on concrete processes, and \emph{distributional} (not merely moment) identities are verified from finite samples. The sampling itself is \emph{reuse} of Units~0/1/2, not new work.
\item[Build] Instantiate each process by reusing the existing routes (Cholesky from Unit~0, KL from Unit~2, and circulant embedding from Unit~1 for the stationary OU); assemble a small kernel catalogue with the relevant closed-form facts attached to each.
\item[Check] (a) \emph{Route equivalence:} Cholesky, KL, and circulant samples of the OU process yield the same empirical covariance to Monte-Carlo tolerance. (b) \emph{Distributional identities:} Brownian-motion invariances tested as distributional identities on empirical covariances (e.g.\ $c^{-1/2}W(ct) \overset{d}{=} W(t)$); bridge endpoints pinned exactly at $t=0,1$ via the covariance/Cholesky route (the time-change formula cannot be evaluated at $t=1$). (c) \emph{The cross-check} (below).
\item[Analytic target] Kernels $\min(t,s)$, $\min(t,s)-ts$, $\tfrac{D}{\alpha}e^{-\alpha|\tau|}$. \textbf{The cross-check target is asymptotic, not exact.} On $[0,T]$ the OU KL eigenvalues do \emph{not} equal the sampled spectral density; the coincidence $\lambda_k=\hat R(k)$ is torus-only. What holds on $[0,T]$ is the large-$T$ Toeplitz/Grenander--Szeg\H{o} limit $\lambda_k \to S(\omega_k)$, $\omega_k = k\pi/T$. \extnote{this is the Toeplitz asymptotic flagged in \S1.2, not a Pavliotis result; it is the reason the cross-check is stated as a convergence, not an equality.}
\item[Pass/fail] Route equivalence: pairwise $\|\widehat\Sigma_{\text{route A}}-\widehat\Sigma_{\text{route B}}\|_F$ within Monte-Carlo tolerance ($\lesssim3\times$ the single-route MC SE). Distributional identity: empirical covariance of $c^{-1/2}W(ct)$ matches $\min(t,s)$ within MC tolerance. Cross-check: the gap $g(T)=\max_k|\lambda_k - S(\omega_k)|$ \emph{decreases monotonically} across a ladder of increasing $T$ (a convergence predicate, not a fixed tolerance).
\item[Negative control] Repeat the self-similarity identity with the \emph{wrong} exponent, $c^{-1/3}W(ct)$: the empirical covariance no longer matches $\min(t,s)$ (it is off by $c^{1/3}$). This makes the correct $c^{-1/2}$ meaningful and is fully self-contained.
\item[Artifact] A three-panel portrait per process (sample paths / covariance heat-map / eigenvalue decay); a route-equivalence panel (three empirical covariances of OU, pairwise differences at MC scale); the cross-check plot $g(T)$ vs.\ $T$ decreasing.
\item[README] Concept: each process is just a choice of kernel. Why this check: distributional identities certify the \emph{law}, not merely a moment, and the cross-check forces the two diagonalizations to reconcile as $T\to\infty$. What it proves: the field's vocabulary is correctly instantiated, the sampling routes agree, and the Fourier and KL vertices are made to speak to each other directly --- the payoff of \S1.5. Its distinct idea is reconciliation; the sampling is inherited.
\item[Missteps] (i) \emph{Expecting exact $\lambda_k=S(\omega_k)$ on $[0,T]$} --- it is asymptotic; gate on the gap \emph{shrinking}, not on a single-$T$ match. (ii) \emph{Frequency grid for the cross-check} --- use $\omega_k=k\pi/T$ and the same $1/2\pi$ convention as Unit~1, or the comparison is off by a constant. (iii) \emph{Distributional vs moment identity} --- equality of mean and covariance implies equality in law only for jointly Gaussian objects \pav{App.~B.5}; do not claim a distributional identity from a single matched moment.
\item[Purpose] The field's basic vocabulary; every subsequent construction is one of these three processes (or a simple transform of one) --- and the one place the whole apparatus is shown to be internally consistent.
\item[Effort] $\approx 7$ h (leverages Units 0, 1, and 2; almost entirely reuse).
\item[Prerequisite] For jointly Gaussian objects, equality of mean and covariance implies equality in law \pav{App.~B.5} --- the recurring fact that uncorrelated $+$ jointly Gaussian $\Rightarrow$ independent is used here to validate KL coefficients. New content: verifying a \emph{distributional} identity (not just a moment) from finite Monte Carlo samples.
\end{unitmeta}

\begin{contracttable}
\texttt{kernels.jl} & \texttt{brownian\_motion\_kernel}, \texttt{brownian\_bridge\_kernel}, \texttt{exponential\_kernel} & the three catalogue kernels \\
\texttt{sampling.jl} & \texttt{sample\_cholesky}, \texttt{sample\_kl}, \texttt{sample\_circulant\_embedding} & reused; route equivalence \\
\texttt{gaussianprocess.jl} & \texttt{empirical\_cov(paths)} & per-route empirical covariance \\
\texttt{kl.jl} / \texttt{spectral.jl} & \texttt{nystrom\_eigen(R, nodes, weights)}, \texttt{welch\_psd(x, dt; ...)} & $\lambda_k$ and $S(\omega_k)$ for the cross-check \\
\end{contracttable}
\noindent\textbf{No new \texttt{src/} code.} \textbf{\texttt{experiments/03\_process\_zoo/run.jl}} produces: (i) three per-process portraits; (ii) the OU route-equivalence panel; (iii) the BM invariance identities and pinned bridge endpoints; (iv) the cross-check ladder $g(T)=\max_k|\lambda_k-S(\omega_k)|$ vs.\ $T$; (v) the wrong-exponent negative-control panel.

% ---------------------------------------------------------------
\subsection*{Unit 4 --- Ergodicity as Loop-Closer}
\addcontentsline{toc}{subsection}{Unit 4 --- Ergodicity as Loop-Closer}
\begin{unitmeta}
\item[Concept] Ergodic properties of stationary processes \pav{\S1.2}; the $L^2$ law of large numbers, Prop.~1.16 \eqref{eq:ergodic}, valid under $C\in L^1(0,\infty)$; the finite-$T$ variance identity, Lemma~1.17 \eqref{eq:ergodic-rate}; the Green--Kubo formula \eqref{eq:greenkubo}; Remark~1.19 (the generalization to $f(X_t)$).
\item[Build] For the stationary OU process, estimate mean and covariance two ways: from \emph{one long path} (time average) and from \emph{many short paths} (ensemble average); track the running time-average as a function of $T$; estimate the variance of the time-average across an ensemble of length-$T$ paths.
\item[Check] The variance of the time-average against the ergodic rate; the Green--Kubo growth of the integrated process; the running average against the ensemble mean.
\item[Analytic target] $\operatorname{Var}\!\left(\frac{1}{T}\int_0^T X_s\,ds\right)\to \frac{2\Dstar}{T} = \frac{2D}{T\alpha^2}$ (both slope and constant, carrying $\alpha$); Green--Kubo: $\mathbb{E}\big(\int_0^t X_s\,ds\big)^2 \to 2\Dstar t = \frac{2D}{\alpha^2}\,t$.
\item[Pass/fail] Fit $\log\operatorname{Var}$ vs.\ $\log T$; gate slope against $-1$ within $2.5\,\mathrm{SE}$ \emph{and} the fitted constant against $2\Dstar$ within $10\%$ (the constant is what pins $\alpha$). Green--Kubo: fitted MSD slope against $2\Dstar$ within $10\%$.
\item[Negative control] Feed the \emph{same} \texttt{time\_average\_variance} a synthetic stationary process with a \emph{non-integrable} correlation, $C(u)\propto(1+u)^{-1/2}$ (so $C\notin L^1$): the variance decays slower than $1/T$ and the fitted slope departs from $-1$. This shows Prop.~1.16's $L^1$ hypothesis is load-bearing --- self-contained, with no forward reference to fBm (Unit~6 does the full fBm version).
\item[Artifact] Running time-average converging to the ensemble mean, plotted with a $\pm\sqrt{2\Dstar/T} = \pm\sqrt{2D/(T\alpha^2)}$ band; variance-vs-$T$ log--log with the fitted $-1$ slope; mean-square-displacement-vs-$t$ with fitted slope $2\Dstar = 2D/\alpha^2$; the non-integrable-$C$ control on the same axes.
\item[README] Concept: ergodicity is the $L^2$ law of large numbers, and correlation decay is what licenses inferring an ensemble fact from a single path. Why this check: the $1/T$ rate and its constant are the quantitative signature, and breaking it on $C\notin L^1$ shows the hypothesis is essential, not decorative. What it proves: the entire project's one-path-vs-analytic methodology (Units 0--3) is justified, and its limits are demonstrated. \emph{This is the methodological keystone --- the strongest single piece of evidence in the repository.}
\item[Missteps] (i) \emph{Confusing the time-average with its variance:} one long path gives a single realization of $\frac1T\int X$, \emph{not} its variance; the variance needs an \emph{ensemble} of length-$T$ paths. Guard: \texttt{time\_average\_variance} takes a path \emph{matrix}. (ii) \emph{Band constant:} compute the $\pm\sqrt{2\Dstar/T}$ band with $\Dstar=D/\alpha^2$, not $D$. (iii) \emph{Integral quadrature:} integrate the path with a consistent trapezoid rule, matching the sampler's grid.
\item[Purpose] The methodological justification, made concrete, for why every theoretical check in Units~0--3 is permitted to compare \emph{one} simulated path against an analytic target.
\item[Effort] $\approx 6$ h.
\item[Prerequisite] Ergodicity is strictly stronger than stationarity --- it additionally requires correlation decay ($C\in L^1$); this $L^2$ law of large numbers is the probabilistic core of the entire project. (Birkhoff's a.s.\ ergodic theorem is a deeper, unnecessary refinement here.) \pav{\S1.2, Prop.~1.16}
\end{unitmeta}

\begin{contracttable}
\texttt{ergodic.jl} & \texttt{running\_time\_average(x, dt)} & cumulative $\frac1t\int_0^t X_s\,ds$ along one path \\
 & \texttt{time\_average\_variance(paths, dt; horizons)} & $(T\text{-values},\ \widehat{\operatorname{Var}})$ across an ensemble \\
 & \texttt{msd\_integrated(paths, dt)} & $(t,\ \mathbb{E}(\int_0^t X)^2)$ \\
 & \texttt{greenkubo\_coefficient(; D, alpha)} & $\Dstar = D/\alpha^2$ \\
\texttt{sampling.jl} & \texttt{sample\_circulant\_embedding} / \texttt{sample\_cholesky} & OU paths (reused) \\
\end{contracttable}
\noindent\textbf{\texttt{experiments/04\_ergodicity/run.jl}} produces: (i) $(T,\ \widehat{\operatorname{Var}})$ with fitted $1/T$ slope and constant vs.\ $2\Dstar$; (ii) $(t,\ \text{MSD})$ with fitted slope vs.\ $2\Dstar$; (iii) the running-average path with the $\pm\sqrt{2\Dstar/T}$ band; (iv) the synthetic non-integrable-$C$ control.

% ---------------------------------------------------------------
\subsection*{Unit 5 --- Brownian Motion as a Scaling Limit (Random Walk $\to$ BM)}
\addcontentsline{toc}{subsection}{Unit 5 --- Brownian Motion as a Scaling Limit}
\begin{unitmeta}
\item[Concept] Brownian motion as the weak limit of a rescaled random walk \pav{\S1.3}: for iid mean-zero, variance-one $\{\xi_j\}$ with $S_n=\sum_{j=1}^n\xi_j$, the interpolated, rescaled walk
\[
W^n_t \;=\; \tfrac{1}{\sqrt{n}}S_{[nt]} + (nt-[nt])\,\tfrac{1}{\sqrt{n}}\xi_{[nt]+1}
\]
converges weakly to a standard Brownian motion as $n\to\infty$. This is BM's \emph{construction}, complementary to the covariance-operator picture: a statement about laws on path space rather than an operation on $\mathcal{C}$.
\item[Build] Simulate rescaled random walks for a ladder of $n$ over \emph{three} increment laws chosen to separate universality of the limit from law-dependence of the rate: \textbf{Rademacher} and \textbf{uniform} (both symmetric) and a \textbf{centered exponential} $\xi = E-1$, $E\sim\mathrm{Exp}(1)$ (asymmetric, skewness $2$); interpolate to a common grid; collect one-time marginals, two-time covariances, and a path functional (running maximum).
\item[Check] Convergence of the finite-dimensional distributions and of a path functional to the Gaussian/BM target; the \emph{rate} of marginal convergence, per law.
\item[Analytic target] One-time marginal $W^n_1 \Rightarrow N(0,1)$; two-time covariance $\to\min(s,t)$; running-maximum distribution converging across laws. \textbf{Marginal-rate target (corrected):} the Kolmogorov--Smirnov statistic decays at a \emph{law-dependent} rate. Berry--Esseen bounds it by $O(n^{-1/2})$, but that rate is achieved only when the skewness is nonzero: for a symmetric increment the leading $n^{-1/2}$ Edgeworth term vanishes, leaving the $O(n^{-1})$ kurtosis term. Hence the target is $\propto n^{-1/2}$ for the \emph{exponential} law and $\propto n^{-1}$ for \emph{Rademacher} and \emph{uniform}. \extnote{Edgeworth expansion / Berry--Esseen; standard, external to Pavliotis. High confidence.}
\item[Pass/fail] Fit $\log(\text{KS})$ vs.\ $\log n$ \emph{per law}; gate the exponential-law slope against $-\tfrac12$ and the two symmetric-law slopes against $-1$, each within $2.5\,\mathrm{SE}$. Covariance: empirical $\to\min(s,t)$ within MC tolerance.
\item[Negative control] Add an \emph{infinite-variance} increment (symmetric Pareto with tail index $<2$): the rescaled walk does not converge to $N(0,1)$ at all --- $W^n_1$ departs from the QQ line, its empirical variance \emph{grows} with $n$ instead of settling at $1$, and $n^{-1/2}$ is simply the wrong normalization. This shows the finite-variance hypothesis of Donsker is load-bearing.
\item[Artifact] Overlaid rescaled-walk paths at increasing $n$; marginal histogram / QQ-plot vs.\ $N(0,1)$; empirical vs.\ $\min(s,t)$ covariance; KS-statistic-vs-$n$ log--log \emph{with one line per law}, the exponential at slope $-\tfrac12$ and the symmetric laws at $-1$; the infinite-variance control on the same QQ axes.
\item[README] Concept: BM is the Donsker scaling limit --- a statement about laws on path space, complementary to the covariance picture. Why this check: the finite-dimensional distributions and a path functional converge, and the marginal \emph{rate} follows Berry--Esseen but depends on skewness, disappearing entirely without finite variance. What it proves: universality of the \emph{limit} but law-dependence of the \emph{rate} --- the subtle probabilistic content of \S1.3, demonstrated rather than asserted.
\item[Missteps] (i) \emph{Expecting $n^{-1/2}$ for the symmetric laws} --- it is $n^{-1}$; a ``too steep'' slope is correct, not a bug. (ii) \emph{Too few replicates per $n$} --- the KS statistic is itself noisy; use enough Monte-Carlo replicates to fit a clean slope. (iii) \emph{Interpolation} --- keep the linear-interpolation term so $W^n_t$ is continuous, or the running-maximum functional is biased.
\item[Purpose] The one unit where a \emph{process} (not merely a moment) is shown to converge, closing the \S1.3 story of where Brownian motion comes from; it also motivates, concretely, why BM is the canonical object the rest of the project revolves around.
\item[Status] \textbf{Committed.}
\item[Effort] $\approx 4$--$5$ h (plus a small buffer for the law-dependent-rate surprise).
\item[Prerequisite] This is weak convergence of processes, which sits slightly \emph{outside} the covariance-operator hub --- a statement about laws on path space, not an operation on $\mathcal{C}$; it is included because BM-as-a-limit is core \S1.3 content. \extnote{Donsker's invariance principle / functional CLT is standard but external to the second-order framing of Chapter~1; Berry--Esseen supplies the $n^{-1/2}$ marginal bound, and the Edgeworth expansion the $n^{-1}$ symmetric-law rate.}
\end{unitmeta}

\begin{contracttable}
\texttt{run.jl}-local & \texttt{rescaled\_walk(xi\_sampler, n, t\_grid, rng)} & one interpolated path $W^n_t$ \\
 & \texttt{ks\_statistic(samples)} & KS distance of $W^n_1$ to $N(0,1)$ \\
\texttt{gaussianprocess.jl} & \texttt{empirical\_cov(paths)} & two-time covariance $\to\min(s,t)$ \\
\end{contracttable}
\noindent\textbf{No new \texttt{src/} code} (weak convergence of laws is not an operation on $\mathcal{C}$; the builders live in the experiment). \textbf{\texttt{experiments/05\_bm\_scaling\_limit/run.jl}} produces: (i) overlaid rescaled paths at increasing $n$; (ii) marginal histogram/QQ vs.\ $N(0,1)$; (iii) $(n,\text{KS})$ log--log, one line per increment law; (iv) empirical vs.\ $\min(s,t)$; (v) running-max across laws; (vi) the infinite-variance control.

% ---------------------------------------------------------------
\subsection*{Unit 6 --- Fractional Brownian Motion}
\addcontentsline{toc}{subsection}{Unit 6 --- Fractional Brownian Motion}
\begin{unitmeta}
\item[Concept] Self-similarity/scaling \pav{eqn.~1.28, Prop.~1.27}; long-range dependence for Hurst parameter $H\neq\tfrac12$; covariance
\begin{equation}
R_H(t,s) = \tfrac{1}{2}\left(|t|^{2H} + |s|^{2H} - |t-s|^{2H}\right).
\label{eq:fbm}
\end{equation}
\pav{Def.~1.26, eqn.~1.27}
\item[Build] Sample via Cholesky (with jitter) and --- more efficiently --- via circulant embedding (Unit~1's Route~4), noting that for fBm the plain symmetric embedding is not always nonnegative definite and a larger / minimal-nonnegative embedding may be required. Reuse \texttt{nystrom\_eigen} for the eigenvalue-decay contrast and \texttt{time\_average\_variance} for the ergodic-breakage control.
\item[Check] Empirical vs.\ analytic covariance for several $H$; the self-similarity identity; the eigenvalue-decay contrast with BM; and the ergodic breakage at $H>\tfrac12$.
\item[Analytic target] $R_H(t,s)=\tfrac12(|t|^{2H}+|s|^{2H}-|t-s|^{2H})$ for $H\in\{0.3,0.5,0.8\}$; self-similarity $W^H_{ct}\overset{d}{=}c^{H}W^H_t$ \pav{eqn.~1.28}; eigenvalue decay of $\mathcal{C}_H$ contrasted with the $k^{-2}$ decay of Brownian motion.
\item[Pass/fail] $\|\widehat R_H-R_H\|_F/\|R_H\|_F <$ MC tolerance for each $H$; self-similarity as a covariance identity within MC tolerance; and --- for the control --- the fitted ergodic slope at $H=0.8$ \emph{departing from} $-1$ (predicate: $|\text{slope}+1| > $ several SE, i.e.\ the $1/T$ rate provably breaks).
\item[Negative control] The full fBm version of Unit~4's control: apply \texttt{time\_average\_variance} to $H=0.8$. The increments have non-summable correlations ($C\notin L^1$), so the $1/T$ ergodic rate breaks --- the first genuine process where Unit~4's machinery fails, run rather than asserted.
\item[Artifact] Sample paths at $H=0.3,0.5,0.8$ (rougher $\to$ smoother); covariance heat-map; self-similarity scaling check; eigenvalue-decay overlay against BM's $k^{-2}$; the $H=0.8$ ergodic-breakage panel.
\item[README] Concept: fBm is a one-exponent generalization of BM tuning both self-similarity and increment correlation. Why this check: the covariance match and self-similarity confirm the kernel fits the hub, while the ergodic breakage at $H>\tfrac12$ demonstrates the boundary of Unit~4's machinery. What it proves: the hub accommodates a qualitatively new kernel, and the ergodic loop-closer's limits are shown, not merely stated.
\item[Missteps] (i) \emph{Circulant non-negativity fails for fBm} --- the plain embedding can have negative eigenvalues; enlarge to the minimal nonnegative (Dietrich--Newsam) embedding. (ii) \emph{Cholesky near-singularity} for $H\to1$ --- report jitter and prefer circulant. (iii) \emph{Variance normalization} --- $R_H(t,t)=|t|^{2H}$; a missing factor $\tfrac12$ is the usual slip.
\item[Purpose] A genuine covariance case study that fits the hub (fBm is just another kernel) yet is qualitatively new: for $H>\tfrac12$ the increment correlations are \emph{non-summable} (long-range dependence), so $C\notin L^1$ and the ergodic loop-closer of Unit~4 does \emph{not} directly apply --- the first process where that machinery breaks, which is precisely why it is instructive.
\item[Status] \textbf{Committed} (promoted from stretch); alongside Unit~7 it is the material to trim first if the buffer is somehow consumed.
\item[Effort] $\approx 5$--$6$ h; flagged for the debugging buffer (circulant non-negativity).
\item[Prerequisite] fBm generalizes BM ($H=\tfrac12$) through a single exponent controlling both self-similarity and increment correlation. \extnote{circulant embedding for fBm (Dietrich--Newsam) and the nonnegative-embedding caveat are standard but external to Pavliotis.}
\end{unitmeta}

\begin{contracttable}
\texttt{kernels.jl} & \texttt{fbm\_kernel(t, s; H=0.5)} & fBm covariance $R_H$ \\
\texttt{sampling.jl} & \texttt{sample\_cholesky} (jittered), \texttt{sample\_circulant\_embedding} & fBm paths (minimal-nonneg caveat) \\
\texttt{kl.jl} & \texttt{nystrom\_eigen(R\_H, nodes, weights)} & eigenvalue-decay contrast vs.\ $k^{-2}$ \\
\texttt{ergodic.jl} & \texttt{time\_average\_variance(paths, dt; horizons)} & the $H=0.8$ breakage control \\
\end{contracttable}
\noindent\textbf{\texttt{experiments/06\_fbm/run.jl}} produces: (i) paths at $H\in\{0.3,0.5,0.8\}$; (ii) empirical-vs-analytic $R_H$ heat-maps; (iii) the self-similarity covariance check; (iv) the $\mathcal{C}_H$ eigenvalue-decay overlay against BM; (v) the ergodic-breakage panel at $H=0.8$.

% ---------------------------------------------------------------
\subsection*{Unit 7 --- SDE Bridge: OU Invariant Law and Strong vs.\ Weak Convergence}
\addcontentsline{toc}{subsection}{Unit 7 --- SDE Bridge: Strong vs.\ Weak Convergence}
\begin{unitmeta}
\item[Concept] The OU stochastic differential equation with \emph{additive} noise,
\begin{equation}
dX_t = -\alpha X_t\,dt + \sqrt{2D}\,dW_t,
\label{eq:ou-sde}
\end{equation}
whose stationary law is $N(0,D/\alpha)$ with covariance $\tfrac{D}{\alpha}e^{-\alpha|\tau|}$ \pav{cf.\ Ch.~2, Example 2.5} --- note $\sqrt{2D}$ makes the stationary variance $\tfrac{(\sqrt{2D})^2}{2\alpha}=D/\alpha$ agree with the Unit~3 kernel; and its Euler--Maruyama (EM) discretization. The strong/weak-convergence \emph{distinction} is demonstrated on the geometric-Brownian-motion SDE with \emph{multiplicative} noise,
\begin{equation}
dY_t = \mu Y_t\,dt + \sigma Y_t\,dW_t, \qquad Y_t = Y_0\exp\!\big((\mu-\tfrac12\sigma^2)t + \sigma W_t\big),
\label{eq:gbm}
\end{equation}
which has a closed-form solution against which both errors can be measured.
\item[Build] An EM integrator (optionally Milstein); apply it to \eqref{eq:ou-sde} (invariant law/covariance vs.\ the exact stationary OU of Unit~3) and to \eqref{eq:gbm} (convergence orders vs.\ the exact solution), sharing the driving path across resolutions via \texttt{coarsen\_increments}.
\item[Check] (i) OU invariant measure; (ii) GBM strong and weak convergence orders; (iii) the two negative-control contrasts below.
\item[Analytic target] OU: empirical law $\to N(0,D/\alpha)$ with covariance $\to\frac{D}{\alpha}e^{-\alpha|\tau|}$. GBM: strong error $\mathbb{E}|Y^{\Delta t}_T-Y_T| \propto \Delta t^{1/2}$ (slope $\tfrac12$); weak error $|\mathbb{E}f(Y_T^{\Delta t})-\mathbb{E}f(Y_T)| \propto \Delta t$ (slope $1$).
\item[Pass/fail] OU: KS distance of the empirical law to $N(0,D/\alpha)$ below tolerance, covariance match within MC tolerance. GBM: fitted strong slope against $\tfrac12$ and weak slope against $1$, each within $2.5\,\mathrm{SE}$, \emph{and} the two fitted slopes separated by more than the sum of their SEs (they must visibly differ).
\item[Negative control] The most striking falsifier in the project: recompute the GBM strong error with \emph{independent} coarse and fine Brownian paths instead of the shared one. The ``strong error'' stops decaying and plateaus at an $O(1)$ constant (the spread between two unrelated solutions) --- showing that strong convergence is a statement about a \emph{shared} noise realization, and converting the ``coupling trap'' misstep into a demonstration. \emph{Near-free contrast (validates the additive-noise correction):} run the same EM integrator on OU (additive) and GBM (multiplicative) --- OU shows strong order $1.0$ (no gap), GBM shows $0.5$ (gap). One integrator, two noise structures, two orders.
\item[Artifact] Log--log plot of GBM strong and weak errors with both slopes annotated (and visibly distinct); the un-coupled strong-error plateau on the same axes; the OU-vs-GBM strong-order contrast panel; a note that deterministic Taylor-error analysis predicts neither slope.
\item[README] Concept: the SDE view is dynamical; EM discretizes it, and strong vs.\ weak convergence are different questions. Why this check: the $0.5$/$1.0$ gap is a multiplicative-noise phenomenon invisible to deterministic Taylor analysis, and --- as the un-coupled control shows --- the strong error is only meaningful when coarse and fine share a driving path. What it proves: the Chapter-3 bridge is understood and the strong/weak distinction (the subtlest SDE-numerics idea) is demonstrated, not asserted.
\item[Missteps] (i) \emph{The coupling trap} --- share $dW$ across resolutions by subsampling the fine path (\texttt{coarsen\_increments}); the un-coupled version is the control, not the method. (ii) \emph{Weak-error MC floor} --- the weak error ($\sim\Delta t$) is swamped by Monte-Carlo error ($\sim M^{-1/2}$) unless $M$ is large enough; the slope flattens exactly as the eigenvalue floor does. Guard: ensure the MC SE is well below the smallest $\Delta t$ bias. (iii) \emph{Choice of $f$} --- use a smooth observable (e.g.\ $f(y)=y$ or $y^2$) for the weak error.
\item[Purpose] The explicit bridge from the covariance/spectral view of Chapter~1 to the dynamical (SDE) view of Chapter~3.
\item[Status] \textbf{Committed}, but the labelled forward pointer to Chapter~3 and therefore the single cut line if the buffer is consumed; nothing in the core depends on it.
\item[Effort] $\approx 8$ h; flagged for the debugging buffer (coupling + weak-error floor).
\item[Prerequisite] $\mathbb{E}|\Delta W| = \sqrt{\Delta t}$, \emph{not} $\Delta t$ --- the reason deterministic error analysis does not transfer. For \emph{multiplicative} noise the Milstein correction is nonzero, so EM strong order drops to $0.5$ while weak order stays $1.0$; for \emph{additive} noise that correction vanishes and EM is already strong order $1.0$, collapsing the gap. \extnote{the OU/GBM SDEs and the Euler--Maruyama/Milstein strong/weak convergence orders are Chapter-3 and standard-SDE-numerics material (e.g.\ Kloeden \& Platen); the additive-noise strong-order-1 result is the reason the strong/weak contrast must be shown on multiplicative noise. Used here only as an explicitly labelled forward pointer, not claimed as Chapter-1 content.}
\end{unitmeta}

\begin{contracttable}
\texttt{sde.jl} & \texttt{euler\_maruyama(x, a, b, dt, dW)} & one EM step \\
 & \texttt{ou\_step(x, dt, dW; D=1.0, alpha=1.0)} & OU (additive) step \\
 & \texttt{gbm\_exact(Y0, mu, sigma, t, Wt)} & GBM closed-form reference \\
 & \texttt{simulate(step, x0, dt, nsteps, dW)} & full path from a step rule \\
 & \texttt{brownian\_path(dt, nsteps, rng)} & \texttt{(W, dW)} driving path + increments \\
 & \texttt{coarsen\_increments(dW\_fine, m)} & \texttt{dW\_coarse} (shares the path across resolutions) \\
\end{contracttable}
\noindent\textbf{\texttt{experiments/07\_sde\_bridge/run.jl}} produces: (i) the OU invariant-law check (histogram/QQ vs.\ $N(0,D/\alpha)$ and covariance vs.\ the exact kernel); (ii) GBM $(\Delta t,\ \text{strong error})$ with the \emph{shared} path and $(\Delta t,\ \text{weak error})$ from many paths, log--log with both slopes; (iii) the un-coupled strong-error plateau; (iv) the OU-vs-GBM strong-order contrast.

% =====================================================================
\section{Repository Structure and Rationale}
% =====================================================================

The repository separates a small, tested \textbf{library} (\texttt{src/}, organized by \emph{operation on $\mathcal{C}$}) from a narrative, reproducible \textbf{experiment gallery} (\texttt{experiments/}, organized by \emph{unit}). This separation is the central design decision.

\begin{lstlisting}
StochProc.jl/
|-- Project.toml, Manifest.toml     # both committed: exact reproducibility
|-- README.md                       # through-line + gallery index
|-- src/
|   |-- StochProc.jl                # module definition, exports
|   |-- kernels.jl                  # R(t,s): min, min-ts, exp, periodic, fBM
|   |-- gaussianprocess.jl          # GP type, grid assembly, empirical cov
|   |-- sampling.jl                 # the 4 square roots: Cholesky / KL /
|   |                               #   random-Fourier / circulant (needs FFTW)
|   |-- kl.jl                       # Nystrom quadrature eigenproblem (symmetrized)
|   |-- spectral.jl                 # Welch periodogram, Bochner transforms
|   |-- ergodic.jl                  # time-average estimators, Green-Kubo
|   `-- sde.jl                      # Euler-Maruyama, Milstein, coupling helpers:
|                                   #   OU (additive) + geometric BM (multiplicative)
|-- test/                           # analytic identities, tight tolerances
|   `-- runtests.jl
|-- experiments/                    # the gallery -- one folder per unit
|   |-- 00_covariance_core/     {run.jl, README.md, figures/}
|   |-- 01_spectral_bochner/    {run.jl, README.md, figures/}
|   |-- 02_kl_quadrature/       {run.jl, README.md, figures/}
|   |-- 03_process_zoo/         {run.jl, README.md, figures/}
|   |-- 04_ergodicity/          {run.jl, README.md, figures/}
|   |-- 05_bm_scaling_limit/    {run.jl, README.md, figures/}
|   |-- 06_fbm/                 {run.jl, README.md, figures/}
|   `-- 07_sde_bridge/          {run.jl, README.md, figures/}   # Ch.3 pointer
`-- .github/workflows/ci.yml
\end{lstlisting}

\textbf{Rationale.} \texttt{src/} is organized by \emph{operation on $\mathcal{C}$} --- kernels, the two diagonalizations, the four square roots, the ergodic loop --- rather than by named process (Brownian motion, OU, \ldots). A process is nothing more than a choice of kernel; organizing by operation makes the module structure a faithful expression of the through-line in Section~1 \emph{for the second-order theory of Chapter~1}. Two units sit deliberately at the edge of that spine and are placed among \texttt{experiments/} accordingly: Unit~5 (random walk $\to$ BM) is weak convergence of \emph{laws}, not an operation on $\mathcal{C}$, and lives entirely in its experiment folder; Unit~7 (SDE bridge) is a dynamical, generator-flavored Chapter-3 pointer whose only library home is \texttt{sde.jl}. A caveat on scope, stated honestly: the covariance-operator spine is the right organizing principle precisely because Chapter~1 is the Gaussian/second-order chapter. Adding later-chapter material never forces a rewrite of existing files --- one adds modules to \texttt{src/} and numbered folders to \texttt{experiments/} --- but the covariance operator ceases to be the hub. Chapter~2 (diffusions, Markov generators, Chapman--Kolmogorov, the Kolmogorov equations) and Chapter~4 (Fokker--Planck) are centered on the \emph{generator}/transition semigroup on generally non-Gaussian processes; they introduce a second, generator-centered spine \emph{alongside} $\mathcal{C}$, rather than more operations on it. So the code accretes without restructuring, but the single-hub \emph{narrative} is Chapter-1-specific and should not be over-sold as the framework's permanent organizing idea. Each experiment folder remains self-contained and documents concept~$\to$~implementation~$\to$~theoretical check~$\to$~result, which directly serves the goal of leaving durable evidence of understanding rather than evidence of having merely read the material.

% =====================================================================
\section{Reproducibility and Testing Convention}
% =====================================================================

\begin{itemize}[leftmargin=1.4cm]
  \item \textbf{Seeding.} Use \texttt{StableRNGs.jl} rather than Julia's default global RNG. \extnote{the default stream is not guaranteed stable across Julia versions, so seeded Monte Carlo checks can silently drift; this is a Julia-ecosystem fact, not a Pavliotis one.} Pass an explicit \texttt{StableRNG(seed)} to every stochastic routine; record the seed in each experiment's README.
  \item \textbf{Numerical positive-definiteness (Cholesky).} Form $\Sigma$ on the grid and factor $\Sigma + \varepsilon I$ with a small jitter (nugget) $\varepsilon$; under grid refinement, smooth kernels (exponential/OU, fBm) produce rapidly decaying eigenvalues and a numerically indefinite $\Sigma$, and naive \texttt{cholesky(Symmetric($\Sigma$))} will throw. Report the jitter used, and prefer the KL or circulant route for the smoothest kernels where even jittered Cholesky loses accuracy.
  \item \textbf{Spectral normalization.} Pin the periodogram and Bochner-FFT normalizations so that the discrete $\int\hat S\,d\omega \to R(0)$, matching the $1/2\pi$-on-$S$ convention of \eqref{eq:spectraldensity}; a mismatched estimator normalization is the most likely source of a spurious factor of $2\pi$ in the Unit-1 check. Document the chosen normalization (and the circulant-embedding scaling) in the relevant experiment README.
  \item \textbf{Two verification tiers.}
  \begin{enumerate}
    \item \texttt{test/}: deterministic analytic identities --- kernel symmetry/PSD, the trace identity \eqref{eq:trace} in its general form $\sum_n\lambda_n = \int_0^T R(t,t)\,dt$ (an assembly/quadrature sanity check, plus the stationary specialization $T\,R(0)$ on a stationary kernel), Brownian-motion eigenvalues vs.\ the closed form \eqref{eq:bm-kl} --- checked with tight absolute tolerance, over the resolved eigenvalue range only (see the Unit~2 numerical caution).
    \item \texttt{experiments/}: Monte Carlo checks with tolerance set to a small multiple of the estimated standard error, seeded so that results are deterministic given the seed. Each experiment README states explicitly which tolerance regime applies.
  \end{enumerate}
  \item \textbf{Rate checks as the primary evidence.} A measured \emph{slope} on a log--log plot is the headline artifact for every unit, and is stronger evidence of correct implementation than a single-point match. Gate the two classes of slope differently:
  \begin{itemize}[leftmargin=1.0cm]
    \item \emph{Stochastic slopes} (Monte Carlo covariance error vs.\ $N^{-1/2}$; Euler--Maruyama strong $0.5$ / weak $1.0$ on GBM; ergodic variance vs.\ $1/T$; the KS-statistic vs.\ $n^{-1/2}$ or $n^{-1}$ of Unit~5): fit the slope on a log--log scale and gate on agreement with the theoretical value to within a small multiple ($2$--$3\times$) of the fitted slope's own standard error --- the estimator variance differs sharply between checks, so a single fixed absolute tolerance would be simultaneously too loose in one place and too tight in another.
    \item \emph{The deterministic KL eigenvalue-decay slope}: here the regression residuals are model misspecification (the asymptotic $k^{-2}$ plus the discretization noise floor), \emph{not} sampling noise, so a regression standard error is not a sampling standard error and the SE-based gate is not statistically meaningful. Fit only within the resolved window (above the noise floor) and gate the fitted exponent against the theoretical $-2$ with a fixed absolute margin.
  \end{itemize}
  \item \textbf{Consolidated risk register (the recurring missteps).} Three failure families recur across units; internalize them once rather than per unit.
  \begin{itemize}[leftmargin=1.0cm]
    \item \emph{The FFT/normalization family} (Units 1, 6, 7): one- vs two-sided spectra; angular ($\omega$) vs ordinary ($f=\omega/2\pi$) frequency, worth a factor $2\pi$; FFTW returns unnormalized transforms (restore $\Delta t$, $1/N$); Welch divides by the window power $\sum_n w_n^2$, not $N$; the circulant-embedding scaling is validated once, in Unit~1. Single guard: pin every normalization so the discrete $\int\widehat S\,d\omega\to R(0)$ and unit-test it.
    \item \emph{The estimator noise floor} (Units 2, 7, and the raw periodogram of Unit 1): every rate check has a level below which the estimator's own variance or discretization dominates and the slope flattens --- the Nystr\"om eigenvalue floor, the weak-error Monte-Carlo floor, and the inconsistent raw periodogram are the \emph{same} phenomenon. Single guard: identify the floor, fit slopes only above it, and state the resolved range.
    \item \emph{The strong-convergence coupling} (Unit 7): strong-error checks require the \emph{same} driving Brownian path at coarse and fine steps --- generate the fine path first and subsample it (\texttt{coarsen\_increments}). The un-coupled version is a negative control, not the method.
  \end{itemize}
  \item \textbf{Continuous integration.} Run \texttt{test/runtests.jl} on push via GitHub Actions --- fast, and it guards the deterministic analytic core of \texttt{src/}. Do \emph{not} run the slower Monte Carlo experiments in CI; run those locally and commit the resulting figures directly.
  \item \textbf{Environment pinning.} Commit \texttt{Manifest.toml} alongside \texttt{Project.toml} so that \texttt{Pkg.instantiate()} reproduces the exact dependency versions used. Manifest pinning together with explicit RNG seeding gives full end-to-end reproducibility of every figure in the gallery.
\end{itemize}

% =====================================================================
\appendix
\section{Repository Skeleton: Illustrative Stubs}
% =====================================================================

The following listings illustrate the intended shape of the core library files and the test harness. These are stubs to be filled in per-unit via detailed specifications, not final implementations.

\subsection*{\texttt{src/kernels.jl}}
\begin{lstlisting}
module Kernels

export brownian_motion_kernel, brownian_bridge_kernel,
       exponential_kernel, periodic_kernel, fbm_kernel

"Covariance R(t,s) = min(t,s): standard Brownian motion on [0,T]."
brownian_motion_kernel(t, s) = min(t, s)

"Covariance R(t,s) = min(t,s) - t*s: Brownian bridge on [0,1]."
brownian_bridge_kernel(t, s) = min(t, s) - t * s

"Stationary exponential kernel R(tau) = (D/alpha) * exp(-alpha*|tau|);
 corresponds to the stationary Ornstein-Uhlenbeck process (Example 1.15).
 Note: D is the noise strength (C(0)=D/alpha); the Green-Kubo transport
 coefficient is Dstar = int_0^inf C = D/alpha^2, distinct unless alpha=1.
 (Example 1.18 computes exactly this: int_0^inf C = C(0)*tau_cor = D/alpha^2.)"
exponential_kernel(t, s; D=1.0, alpha=1.0) = (D / alpha) * exp(-alpha * abs(t - s))

"Stationary periodic kernel on the torus [0,1), exp of geodesic distance;
 used for the KL/Bochner coincidence contrast. Positive-definite on the
 circle: its Fourier coefficients Rhat(k) = 2*alpha/(alpha^2+(2*pi*k)^2)
 * (1 - (-1)^k * exp(-alpha/2)) are all > 0 (derived and verified in-house)."
periodic_kernel(t, s; D=1.0, alpha=1.0) =
    (D / alpha) * exp(-alpha * min(abs(t - s), 1 - abs(t - s)))

"Covariance R_H(t,s) = 0.5*(|t|^{2H} + |s|^{2H} - |t-s|^{2H}):
 fractional Brownian motion with Hurst parameter H in (0,1).
 H = 1/2 recovers standard Brownian motion; H > 1/2 gives long-range
 dependence (non-summable increment correlations, so C not in L^1)."
fbm_kernel(t, s; H=0.5) =
    0.5 * (abs(t)^(2H) + abs(s)^(2H) - abs(t - s)^(2H))

end # module
\end{lstlisting}

\subsection*{\texttt{src/gaussianprocess.jl}}
\begin{lstlisting}
module GaussianProcesses

using LinearAlgebra
export GaussianProcess, assemble_cov, assemble_mean, empirical_cov

"A second-order Gaussian process over a covariance kernel R(t,s) and an
 optional mean function. A process is nothing more than a choice of kernel;
 mean and covariance determine the law (Appendix B.5)."
struct GaussianProcess{K,M}
    kernel::K          # R(t, s)
    meanfn::M          # m(t)
end
GaussianProcess(kernel; meanfn = t -> 0.0) = GaussianProcess(kernel, meanfn)

"Assemble Sigma[i,j] = R(t_i, t_j) on a grid. Returned Symmetric so the
 downstream Cholesky/eigen see an exactly symmetric operand (floating-point
 assembly is not bitwise symmetric otherwise)."
function assemble_cov(gp::GaussianProcess, t_grid::AbstractVector)
    R = gp.kernel
    Sigma = [R(t, s) for t in t_grid, s in t_grid]
    return Symmetric(Sigma)
end

"Assemble the mean vector m(t_i) on a grid."
assemble_mean(gp::GaussianProcess, t_grid::AbstractVector) =
    [gp.meanfn(t) for t in t_grid]

"Empirical covariance across sample paths. `paths` is n_grid x N (one path
 per COLUMN -- a transpose silently estimates the wrong matrix). The
 Frobenius error vs. assemble_cov decays at the Monte Carlo rate N^{-1/2}
 (the Unit-0 headline check)."
function empirical_cov(paths::AbstractMatrix)
    n, N = size(paths)
    mu = sum(paths, dims = 2) ./ N
    Xc = paths .- mu
    return (Xc * Xc') ./ (N - 1)
end

end # module
\end{lstlisting}

\subsection*{\texttt{src/sampling.jl}}
\begin{lstlisting}
module Sampling

using LinearAlgebra, FFTW
export sample_cholesky, sample_kl, sample_random_fourier,
       sample_circulant_embedding

"Route 1 (general): (jittered) Cholesky factorization of the covariance matrix.
 Sigma = L * L', X = L * z, z ~ N(0, I). The nugget eps*I restores
 numerical positive-definiteness for smooth, rapidly-decaying kernels."
function sample_cholesky(Sigma::AbstractMatrix, rng; jitter=1e-10)
    n = size(Sigma, 1)
    L = cholesky(Symmetric(Sigma) + jitter * I(n)).L
    return L * randn(rng, n)
end

"Route 2 (general): truncated Karhunen-Loeve expansion.
 X_t = sum_k sqrt(lambda_k) * xi_k * e_k(t), xi_k ~ iid N(0,1)."
function sample_kl(lambdas::AbstractVector, eigenfuncs::AbstractMatrix, rng)
    K = length(lambdas)
    xi = randn(rng, K)
    return eigenfuncs * (sqrt.(lambdas) .* xi)
end

"Route 3 (stationary only): random Fourier series with GAUSSIAN amplitudes,
 on a ONE-SIDED frequency grid omega_grid (omega_k > 0; exclude 0):
 X(t) = sum_k [a_k cos(omega_k t) + b_k sin(omega_k t)],
 a_k, b_k ~ iid N(0, 2 * S(omega_k) * dOmega).
 The factor 2 folds the two-sided spectrum onto omega > 0, so that
 Cov(tau) = sum_k 2*S(omega_k)*dOmega * cos(omega_k*tau) -> R(tau).
 Gaussian (not fixed-amplitude random-phase) amplitudes make this a genuine
 square root of C: the finite-K sample is EXACTLY Gaussian and reproduces the
 law, not merely the covariance. (A random-phase foil with fixed amplitude
 sqrt(4*S(omega_k)*dOmega) and uniform phase matches the covariance but is
 Gaussian only in the K -> infinity limit, so it is not a square root.)"
function sample_random_fourier(S::Function, t_grid, omega_grid, rng)
    dOmega = step(omega_grid)          # omega_grid one-sided: 0 < omega_1 < ...
    X = zeros(length(t_grid))
    for w in omega_grid
        sigma = sqrt(2 * max(S(w), 0.0) * dOmega)
        a = sigma * randn(rng)
        b = sigma * randn(rng)
        @. X += a * cos(w * t_grid) + b * sin(w * t_grid)
    end
    return X
end

"Route 4 (stationary only): circulant (FFT) embedding -- exact O(m log m)
 sampling on a uniform grid (Wood-Chan / Dietrich-Newsam). r = [R(0), R(dt),
 ..., R((n-1)*dt)] is the first column of the Toeplitz covariance. Embed into
 a symmetric circulant of size m = 2(n-1); its eigenvalues are lambda =
 real(fft(c)). Requires lambda .>= 0 (holds for the exponential/periodic
 kernels; for fBm a larger / minimal-nonnegative embedding may be needed).
 This is the Bochner square root made computational: FFT-diagonalization of
 the (circulant) covariance. NOTE: the exact scaling is the per-unit spec
 detail -- validate against the analytic covariance in Unit 1."
function sample_circulant_embedding(r::AbstractVector, rng)
    n = length(r)
    c = vcat(r, r[end-1:-1:2])                 # length m = 2(n-1)
    m = length(c)
    lambda = real(fft(c))                      # circulant eigenvalues
    all(lambda .>= -1e-10) || error("circulant embedding not PSD; enlarge m")
    xi = randn(rng, m) .+ im .* randn(rng, m)
    Y  = fft(sqrt.(max.(lambda, 0.0)) .* xi) ./ sqrt(m)
    return real(Y[1:n])                        # imag(Y[1:n]) is a 2nd draw
end

end # module
\end{lstlisting}

\subsection*{\texttt{src/kl.jl}}
\begin{lstlisting}
module KL

using LinearAlgebra
export quad_nodes_weights, nystrom_eigen, trace_diag, kl_tail_energy

"Quadrature nodes and weights on [0, T]. Trapezoid by default; any positive
 rule works (the Nystrom eigenproblem only needs w_i > 0)."
function quad_nodes_weights(T; n, rule = :trapezoid)
    nodes = range(0, T; length = n)
    if rule == :trapezoid
        w = fill(step(nodes), n)
        w[1] *= 0.5
        w[end] *= 0.5
    else
        error("unknown rule: $rule")
    end
    return collect(nodes), w
end

"Nystrom discretization of  int_0^T R(t,s) e(s) ds = lambda * e(t).
 CRITICAL (symmetrization): the naive discretization (K*W) e = lambda e is
 NON-symmetric even though R is symmetric (W = Diagonal(weights)), so its
 eigenvectors are not orthonormal and can come out complex. Solve the
 SYMMETRIC problem (D*K*D) g = lambda g with D = Diagonal(sqrt.(w)), then
 recover e = D^{-1} g. Eigenpairs are returned sorted DESCENDING and
 L2-normalized on the grid (sum_i w_i e_k(t_i)^2 = 1)."
function nystrom_eigen(R, nodes, weights; nev = length(nodes))
    K = [R(t, s) for t in nodes, s in nodes]
    d = sqrt.(weights)
    A = Symmetric((d * d') .* K)               # D*K*D with D = Diagonal(d)
    vals, vecs = eigen(A)                       # ascending
    idx = sortperm(vals; rev = true)[1:nev]
    lambdas = vals[idx]
    efns = vecs[:, idx] ./ d                    # e = D^{-1} g (columns)
    return lambdas, efns
end

"Grid quadrature of the diagonal:  sum_i w_i R(t_i, t_i) ~ int_0^T R(t,t) dt.
 This is tr(C) in the GENERAL (possibly non-stationary) form and is the
 assembly/quadrature sanity check of Unit 2 -- it validates the weights and
 matrix assembly, NOT the eigensolver. For a stationary kernel R(t,t)=R(0)
 it collapses to T*R(0) (Exercise 23)."
trace_diag(R, nodes, weights) =
    sum(w * R(t, t) for (t, w) in zip(nodes, weights))

"Relative KL truncation (tail) energy beyond K modes:
 sum_{k>K} lambda_k / sum_k lambda_k.  Drives the Exercise-28 cost study."
kl_tail_energy(lambdas::AbstractVector, K::Integer) =
    sum(lambdas[K+1:end]) / sum(lambdas)

end # module
\end{lstlisting}

\subsection*{\texttt{src/spectral.jl}}
\begin{lstlisting}
module Spectral

using FFTW
export welch_psd, raw_periodogram, bochner_forward, bochner_inverse,
       spectral_variance

# Fold a two-sided (omega, S) pair onto omega >= 0, doubling interior bins so
# that the one-sided integral still returns int S dOmega over the full line.
function _onesided(omega, S)
    keep = omega .>= 0
    o = omega[keep]; s = copy(S[keep])
    s[2:end] .*= 2                               # double all but the DC bin
    return o, s
end

"Welch (averaged, windowed) PSD estimate of a real series x sampled at dt.
 CONSISTENT: its variance shrinks as the number of segments grows. Normalized
 so the discrete int Shat dOmega -> R(0) under the 1/2pi-on-S convention of
 (1.7)-(1.8): angular frequency omega = 2*pi*f, and the window power
 U = sum(win.^2) is divided out (NOT the segment length). Returns a ONE-SIDED
 spectrum by default. NOTE: the exact fold/window convention is the Unit-1
 spec detail -- pin it so int Shat dOmega -> R(0) and unit-test that gate."
function welch_psd(x, dt; nseg, noverlap = 0, window = :hann, onesided = true)
    N = length(x); L = div(N, nseg)
    win = window === :hann ?
          [0.5 - 0.5 * cos(2pi * k / (L - 1)) for k in 0:L-1] : ones(L)
    U = sum(abs2, win)                            # window power
    hop = L - noverlap
    acc = zeros(L)
    nused = 0
    start = 1
    while start + L - 1 <= N
        seg = @view x[start:start+L-1]
        acc .+= abs2.(fft(win .* seg))
        nused += 1; start += hop
    end
    Sfull = (dt / (2pi * U)) .* (acc ./ nused)    # 1/2pi-on-S normalization
    omega_full = 2pi .* fftfreq(L, 1 / dt)
    p = sortperm(omega_full)
    omega_full, Sfull = omega_full[p], Sfull[p]
    return onesided ? _onesided(omega_full, Sfull) : (omega_full, Sfull)
end

"Raw (single-shot) periodogram -- the INCONSISTENT estimator kept ONLY as the
 Unit-1 negative control: its variance does NOT decay with record length at
 fixed omega. Same normalization as welch_psd so the two overlay honestly."
raw_periodogram(x, dt; onesided = true) =
    welch_psd(x, dt; nseg = 1, window = :none, onesided = onesided)

"Bochner forward transform of a covariance sequence r = [R(0), R(dt), ...]
 into S(omega): the discrete realization of
 S(omega) = (1/2pi) int e^{-i omega t} R(t) dt   (1.7)."
function bochner_forward(r, dt; onesided = true)
    rsym = vcat(r, r[end-1:-1:2])                 # even extension
    S = real(fft(rsym)) .* (dt / (2pi))
    omega = 2pi .* fftfreq(length(rsym), 1 / dt)
    p = sortperm(omega); omega, S = omega[p], S[p]
    return onesided ? _onesided(omega, S) : (omega, S)
end

"Bochner inverse: S(omega) -> covariance, the discrete form of
 R(t) = int e^{i omega t} S(omega) dOmega   (1.8)."
function bochner_inverse(S, dOmega)
    return real(ifft(S)) .* (length(S) * dOmega)
end

"Total-variance check: trapezoidal int Shat dOmega, which must approach R(0).
 A result of 2*pi*R(0) instead of R(0) is the signature of a dropped 1/2pi."
spectral_variance(omega, Shat) =
    sum(0.5 .* (Shat[1:end-1] .+ Shat[2:end]) .* diff(omega))

end # module
\end{lstlisting}

\subsection*{\texttt{src/ergodic.jl}}
\begin{lstlisting}
module Ergodic

using Statistics
export running_time_average, time_average_variance, msd_integrated,
       greenkubo_coefficient

"Running time-average (1/t) int_0^t X_s ds along ONE path, by cumulative
 trapezoidal quadrature on a uniform grid of spacing dt. This is the object
 whose convergence to the mean is the Unit-4 loop-closer."
function running_time_average(x, dt)
    n = length(x)
    integral = zeros(n)
    acc = 0.0
    for i in 2:n
        acc += 0.5 * (x[i] + x[i-1]) * dt          # trapezoid increment
        integral[i] = acc
    end
    t = dt .* (0:n-1)
    avg = similar(integral)
    avg[1] = x[1]                                   # t -> 0 limit
    @. avg[2:end] = integral[2:end] / t[2:end]
    return avg
end

"Empirical Var((1/T) int_0^T X_s ds) as a function of horizon T, estimated
 across an ENSEMBLE of length-T paths. One long path gives a single
 realization of the time-average, NOT its variance -- the common trap; hence
 `paths` is a MATRIX (n_grid x M). Returns (T-values, variance estimates)."
function time_average_variance(paths, dt; horizons)
    M = size(paths, 2)
    Tvals = Float64[]; varhat = Float64[]
    for h in horizons
        k = round(Int, h / dt)                     # grid index of the horizon
        avgs = [running_time_average(view(paths, 1:k, m), dt)[end] for m in 1:M]
        push!(Tvals, (k - 1) * dt)
        push!(varhat, var(avgs))
    end
    return Tvals, varhat
end

"Mean-square displacement of the integral E( int_0^t X_s ds )^2 across paths;
 for a stationary process it grows ~ 2*Dstar*t (Green-Kubo, (1.14)-(1.15))."
function msd_integrated(paths, dt)
    n, M = size(paths)
    t = dt .* (0:n-1)
    msd = zeros(n)
    for m in 1:M
        acc = 0.0
        I = zeros(n)
        for i in 2:n
            acc += 0.5 * (paths[i, m] + paths[i-1, m]) * dt
            I[i] = acc
        end
        @. msd += I^2
    end
    return t, msd ./ M
end

"Green-Kubo / transport coefficient for the OU kernel:
 Dstar = int_0^inf C = D/alpha^2 (Example 1.18). Distinct from the noise
 strength D unless alpha = 1."
greenkubo_coefficient(; D = 1.0, alpha = 1.0) = D / alpha^2

end # module
\end{lstlisting}

\subsection*{\texttt{src/sde.jl} (excerpt)}
\begin{lstlisting}
module SDE

export euler_maruyama, ou_step, gbm_exact,
       simulate, brownian_path, coarsen_increments

"One Euler-Maruyama step for dX = a(X) dt + b(X) dW. Sharing dW across
 coarse/fine resolutions is the caller's responsibility (strong-error trap)."
euler_maruyama(x, a, b, dt, dW) = x + a(x) * dt + b(x) * dW

"OU with ADDITIVE noise: dX = -alpha X dt + sqrt(2D) dW. Stationary law
 N(0, D/alpha), covariance (D/alpha) exp(-alpha|tau|). Additive noise =>
 Euler-Maruyama already has STRONG order 1.0 (Milstein term vanishes), so
 strong and weak orders coincide -- the strong/weak GAP is NOT visible here."
ou_step(x, dt, dW; D=1.0, alpha=1.0) =
    euler_maruyama(x, u -> -alpha * u, _ -> sqrt(2D), dt, dW)

"Geometric BM with MULTIPLICATIVE noise: dY = mu Y dt + sigma Y dW.
 Exact solution used as the reference for both strong and weak error;
 multiplicative noise => EM strong order 0.5, weak order 1.0 (the GAP)."
gbm_exact(Y0, mu, sigma, t, Wt) = Y0 * exp((mu - sigma^2 / 2) * t + sigma * Wt)

"Integrate a path from a step rule and a PRECOMPUTED increment vector dW
 (length nsteps). Passing dW in explicitly is what lets the caller share one
 driving path across resolutions."
function simulate(step, x0, dt, nsteps, dW::AbstractVector)
    x = x0; path = Vector{typeof(x0)}(undef, nsteps + 1); path[1] = x0
    for i in 1:nsteps
        x = step(x, dt, dW[i]); path[i+1] = x
    end
    return path
end

"A driving Brownian path and its increments on a uniform grid:
 dW ~ N(0, dt), W = cumulative sum. Returns (W, dW)."
function brownian_path(dt, nsteps, rng)
    dW = sqrt(dt) .* randn(rng, nsteps)
    return vcat(0.0, cumsum(dW)), dW
end

"Coarsen fine increments to a step m times larger by summing each block of m
 (Brownian increments are additive). This is the coupling that makes the
 STRONG error meaningful: the coarse path is driven by the SAME noise as the
 fine one. Using independent draws instead is the Unit-7 negative control."
coarsen_increments(dW_fine::AbstractVector, m::Integer) =
    [sum(@view dW_fine[(j-1)*m+1 : j*m]) for j in 1:div(length(dW_fine), m)]

end # module
\end{lstlisting}

\subsection*{\texttt{test/runtests.jl}}
\begin{lstlisting}
using Test, StochProc, StableRNGs

# Assume a grid t_grid, kernels R (nonstationary) and R_stat (stationary),
# horizon T, grid size n_grid, and an assembled Sigma are defined above.

@testset "Covariance operator basics" begin
    # symmetry and positive semi-definiteness of the kernel catalogue
    @test issymmetric(Sigma)
    @test all(eigvals(Sigma) .>= -1e-10)
end

@testset "KL trace identity (general): sum(lambda) == tr(C) == int_0^T R(t,t) dt" begin
    # general form holds for ANY second-order process, stationary or not.
    # This validates quadrature weights + assembly, NOT the eigensolver.
    nodes, w = quad_nodes_weights(T; n = n_grid)
    lambdas, _ = nystrom_eigen(R, nodes, w)
    trC = trace_diag(R, nodes, w)         # int_0^T R(t,t) dt via the quadrature
    @test isapprox(sum(lambdas), trC; rtol = 1e-3)
end

@testset "Trace: stationary specialization (Exercise 23), stationary kernel only" begin
    # For a STATIONARY kernel R(t,t) == R(0), so tr(C) == T*R(0).
    # (This does NOT hold for Brownian motion, where R(0,0)=0 but tr(C)=T^2/2.)
    nodes, w = quad_nodes_weights(T; n = n_grid)
    lambdas_stat, _ = nystrom_eigen(R_stat, nodes, w)
    @test isapprox(sum(lambdas_stat), T * R_stat(0.0, 0.0); rtol = 1e-3)
end

@testset "Brownian motion KL closed form (Example 1.29), resolved range" begin
    n = 1:10
    lambda_analytic = T^2 .* (2 ./ ((2n .- 1) .* pi)).^2   # [0,T]: lambda_n prop T^2
    @test isapprox(lambdas_numeric[1:10], lambda_analytic; rtol = 1e-2)
    # sanity: sum of analytic eigenvalues == T^2/2 == BM trace (not T*R(0)=0).
    @test isapprox(sum(lambdas_numeric), T^2 / 2; rtol = 1e-2)
end
\end{lstlisting}

% =====================================================================
\section*{Sources}
\addcontentsline{toc}{section}{Sources}
% =====================================================================
\begin{itemize}[leftmargin=1.4cm]
  \item G.~A.~Pavliotis, \emph{Stochastic Processes and Applications}. The project's copy is the \emph{November 11, 2015 draft}; all section, theorem, and equation numbers cited in this document (Chapter~1, plus \S B.5 and a single Chapter-2 pointer, Example~2.5) refer to that draft and were checked directly against it. Note that its numbering differs from the published Springer edition. In particular, \textbf{Mercer's theorem is stated inline and unnumbered in \S1.5} of this draft (there is no ``Theorem~1.18''; ``Example~1.18'' is the exponential-covariance / diffusion-coefficient example, where $\int_0^\infty C = D/\alpha^2$ is computed explicitly), and the ergodic result appears as \emph{Proposition~1.16} in the statement (its proof header inconsistently reads ``Theorem~1.16'').
  \item External references, flagged inline where used and not verified against Pavliotis: Grenander \& Szeg\H{o}, \emph{Toeplitz Forms and Their Applications} (asymptotic KL/spectral-density coincidence off the torus); D.~B.~Percival \& A.~T.~Walden, \emph{Spectral Analysis for Physical Applications}, or P.~J.~Brockwell \& R.~A.~Davis, \emph{Time Series: Theory and Methods} (periodogram inconsistency); C.~R.~Dietrich \& G.~N.~Newsam, and A.~T.~A.~Wood \& G.~Chan (circulant embedding of stationary Gaussian fields); P.~E.~Kloeden \& E.~Platen, \emph{Numerical Solution of Stochastic Differential Equations} (Euler--Maruyama/Milstein strong and weak convergence orders, including strong order~1 for additive noise); W.~Feller, \emph{An Introduction to Probability Theory and Its Applications, Vol.~II} (Edgeworth expansion / the skewness-dependent Berry--Esseen rate of Unit~5); I.~Karatzas \& S.~Shreve, \emph{Brownian Motion and Stochastic Calculus}, \S3.3 (white noise as a generalized process); and the functional central limit theorem / Donsker's invariance principle with Berry--Esseen (Unit~5).
\end{itemize}

\end{document}
